\documentclass[11pt,respuestas,a4]{aleph-examen}
% \documentclass[11pt,a5]{aleph-examen}

% -- Paquetes extra
\usepackage{aleph-comandos}
\usepackage{booktabs}
\usepackage{multicol}
\usepackage{systeme}

% -- Datos
\institucion{Facultad de Ciencias Exactas, Naturales y Ambientales}
\carrera{Catalogo STEM}
\asignatura{Algebra lineal}
\tema{Examen no. 2: Espacios Vectoriales}
\autor{Andres Merino}
\fecha{Semestre 2024-1}

\logouno[0.16\textwidth]{Logos/logoPUCE_04_ac}
\definecolor{colortext}{HTML}{1957d1}
\definecolor{colordef}{HTML}{1957d1}
\fuente{montserrat}

\begin{document}

\encabezado

\section*{Ejercicios}

\begin{preguntas}

%%%%%%%%%%%%%%%%%%%%%%%%%%%%%%%%%%%%%%%%
%% Conjunto generado
%%%%%%%%%%%%%%%%%%%%%%%%%%%%%%%%%%%%%%%%
\item
    En $\R^4$, determine el conjunto generado por:
    \[
        S = \big\{
            (1,0,1,0),\
            (2,1,2,0),\
            (1,1,1,1)\big\}.
    \]
    Escriba el conjunto en funcion de sus restricciones.\puntaje{10}

\begin{respuesta}
    Tomemos $(w,x,y,z)\in\R^4$ tal que
    \begin{align*}
        (w,x,y,z)
        & = a(1,0,1,0) + b(2,1,2,0) + c(1,1,1,1)\\
        & = (a+2b+c,\ b+c,\ a+2b+c,\ c).
    \end{align*}
    Obtenemos el sistema
    \[
        \systeme[abc]{
            a+2b+c = w{,},
            b+c = x{,},
            a+2b+c = y{,},
            c = z{.}
        }
    \]
    Planteamos la matriz ampliada y realizamos reduccion por filas:
    \[
        \begin{pmatrix}
            1 & 2 & 1 & | & w\\
            0 & 1 & 1 & | & x\\
            1 & 2 & 1 & | & y\\
            0 & 0 & 1 & | & z
        \end{pmatrix}
        \sim
        \begin{pmatrix}
            1 & 2 & 1 & | & w\\
            0 & 1 & 1 & | & x\\
            0 & 0 & 1 & | & z\\
            0 & 0 & 0 & | & - w + y
        \end{pmatrix}.
    \]
    Este sistema tiene solucion si y solo si $-w + y = 0$. Por lo tanto, el conjunto generado por $S$ es
    \[
        \{
            (w,x,y,z)\in\R^4 : -w + y = 0
        \}.\qedhere
    \]
\end{respuesta}

%%%%%%%%%%%%%%%%%%%%%%%%%%%%%%%%%%%%%%%%
%% Independencia lineal
%%%%%%%%%%%%%%%%%%%%%%%%%%%%%%%%%%%%%%%%
\item
    En $\R^3$, para $\alpha \in \R$, se toma el conjunto
    \[
        S = \big\{
            (1,0,1),\
            (2\alpha,1,2),\
            (1,\alpha,1)\big\}.
    \]
    Para que valores de $\alpha$ se tiene que $S$ es un conjunto linealmente independiente?\puntaje{10}

\begin{respuesta}
    Tomemos $a,b,c\in\R$ tales que
    \[
        a(1,0,1) + b(2\alpha,1,2) + c(1,\alpha,1) = (0,0,0),
    \]
    operando, se tiene que
    \[
        (a + 2\alpha b + c, b + \alpha c, a + 2b + c) = (0,0,0).
    \]
    Con esto, se obtiene el sistema
    \[
        \systeme[abc]{
            a + 2\alpha b + c = 0{,},
            b + \alpha c = 0{,},
            a + 2b + c = 0{.}
        }
    \]
    Para que el conjunto sea linealmente independiente se necesita que el sistema tenga solucion unica. Para esto, planteamos la matriz ampliada y realizamos reduccion por filas:
    \[
        \begin{pmatrix}
            1 & 2\alpha & 1 & | & 0\\
            0 & 1 & \alpha & | & 0\\
            1 & 2 & 1 & | & 0
        \end{pmatrix}
        \sim
        \begin{pmatrix}
            1 & 2\alpha & 1 & | & 0\\
            0 & 1 & \alpha & | & 0\\
            0 & 0 & 2\alpha^2 -2\alpha & | & 0
        \end{pmatrix}.
    \]
    Tenemos que el sistema tiene solucion unica si y solo si
    \[
        2\alpha^2 - 2\alpha \neq 0.
    \]
    Es decir, esto ocurre cuando $\alpha\neq 0$ y $\alpha\neq 1$.
    Con esto, se obtiene que $S$ es linealmente independiente para todo $\alpha$ en $\R\setminus\{0,1\}$.\qedhere
\end{respuesta}

%%%%%%%%%%%%%%%%%%%%%%%%%%%%%%%%%%%%%%%%
%% Subespacios vectoriales y base
%%%%%%%%%%%%%%%%%%%%%%%%%%%%%%%%%%%%%%%%
\item
    En $\R^{2\times2}$, se define el conjunto
    \[
        F = \left\{ \begin{pmatrix} a & b\\ c & d \end{pmatrix} : a-b+c=0\ \land\ b+2d=0 \right\}.
    \]
    \begin{enumerate}
    \item
        Determine una base para $F$.\puntaje{12}
    \item
        Cual es la dimension de $F$?\puntaje{3}
    \end{enumerate}

\begin{respuesta}\hspace{0pt}
    \begin{enumerate}
    \item
        Tomemos un elemento arbitrario de $F$:
        \[
            \begin{pmatrix}
                a & b\\
                c & d
            \end{pmatrix}.
        \]
        Por la definicion de $F$, se cumple
        \[
            a-b+c=0
            \texty
            b+2d=0.
        \]
        De la segunda ecuacion se obtiene $b=-2d$. Sustituyendo en la primera:
        \[
            a-(-2d)+c=0,
        \]
        de donde
        \[
            c=-a-2d.
        \]
        Por lo tanto,
        \[
            \begin{pmatrix}
                a & b\\
                c & d
            \end{pmatrix}
            =
            \begin{pmatrix}
                a & -2d\\
                -a-2d & d
            \end{pmatrix}
            =
            a\begin{pmatrix}
                1 & 0\\
                -1 & 0
            \end{pmatrix}
            +
            d\begin{pmatrix}
                0 & -2\\
                -2 & 1
            \end{pmatrix}.
        \]
        De esta manera, una posible base para $F$ esta dada por
        \[
            \left\{
            \begin{pmatrix}
                1 & 0\\
                -1 & 0
            \end{pmatrix},
            \begin{pmatrix}
                0 & -2\\
                -2 & 1
            \end{pmatrix}
            \right\}.
        \]
        En efecto, este conjunto genera a $F$. Ademas, si
        \[
            \alpha_1\begin{pmatrix}
                1 & 0\\
                -1 & 0
            \end{pmatrix}
            +
            \alpha_2\begin{pmatrix}
                0 & -2\\
                -2 & 1
            \end{pmatrix}
            =
            \begin{pmatrix}
                0 & 0\\
                0 & 0
            \end{pmatrix},
        \]
        al comparar las entradas $(1,1)$ y $(2,2)$ se obtiene $\alpha_1=0$ y $\alpha_2=0$. Por lo tanto, el conjunto es linealmente independiente.
    \item
        Como la base encontrada tiene dos elementos, se concluye que la dimension de $F$ es $2$. \qedhere
    \end{enumerate}
\end{respuesta}

%%%%%%%%%%%%%%%%%%%%%%%%%%%%%%%%%%%%%%%%
%% Aplicaciones lineales
%%%%%%%%%%%%%%%%%%%%%%%%%%%%%%%%%%%%%%%%
\item
    Sea $T: \mathbb{R}^3 \to \mathbb{R}^2$ una aplicacion dada por $T(x,y,z) = (x+2y,\ 3x+4z)$.
    \begin{enumerate}
    \item
        Sin realizar ningun calculo, indique si esta aplicacion podria ser inyectiva o si podria ser sobreyectiva.\puntaje{3}
    \item
        Calcule el nucleo de la aplicacion. \puntaje{5}
    \item
        Determine la imagen de la aplicacion. \puntaje{5}
    \item
        Con lo calculado en los puntos anteriores conteste: La aplicacion es inyectiva? Es sobreyectiva? Es isomorfismo? \puntaje{2}
    \end{enumerate}

\begin{respuesta}
\begin{enumerate}[leftmargin=*]
\item
    Como $\dim(\R^3)=3$ y $\dim(\R^2)=2$, una aplicacion lineal de $\R^3$ en $\R^2$ no puede ser inyectiva. Por otro lado, si podria ser sobreyectiva.
\item
    Tomemos un vector $(x,y,z)$ tal que $T(x,y,z)=(0,0)$. Entonces:
    \[
        T(x,y,z)=(x+2y,\ 3x+4z)=(0,0).
    \]
    Esto produce el sistema
    \[
        \systeme[xyz]{
            x+2y = 0,
            3x+4z = 0.
        }
    \]
    De la primera ecuacion, $x=-2y$. Sustituyendo en la segunda:
    \[
        3(-2y)+4z=0,
    \]
    por lo que
    \[
        z=\frac{3}{2}y.
    \]
    Tomando $y=2t$, se obtiene
    \[
        (x,y,z)=(-4t,2t,3t)=t(-4,2,3).
    \]
    Por lo tanto,
    \[
        \ker(T)=\{t(-4,2,3):t\in\R\}.
    \]
\item
    La imagen esta generada por las imagenes de la base canonica de $\R^3$:
    \[
        T(1,0,0)=(1,3),\qquad
        T(0,1,0)=(2,0),\qquad
        T(0,0,1)=(0,4).
    \]
    Como
    \[
        (2,0)
        \quad\text{y}\quad
        (0,4)
    \]
    son linealmente independientes en $\R^2$, generan todo $\R^2$. Por lo tanto,
    \[
        \img(T)=\R^2.
    \]
\item
    Como el nucleo de la aplicacion no es igual a $\{(0,0,0)\}$, la aplicacion no es inyectiva. Por otro lado, como $\img(T)=\R^2$, la aplicacion es sobreyectiva. Con esto, tenemos que la aplicacion no es biyectiva, por lo que no es un isomorfismo.\qedhere
\end{enumerate}
\end{respuesta}

\end{preguntas}

\end{document}
